\documentclass[exercice]{thedocument}%
\usepackage{texmaths-boxes}%
\startdatakeys
\source{}
\cycle{}
\competencesDuSocle{}
\connaissancesRequises{}
\competencesTravaillees{}
\enddatakeys
\begin{document}%
\begin{enumerate}
    \item Voici un tableau de valeurs d'une fonction $f$ :\par
    \begin{tblr}{colspec=*{8}{Q[c,m]},hlines,vlines}
        $x$&-2&-1&0&1&3&4&5\\
        $f(x)$&5&3&1&-1&-5&-7&-9
    \end{tblr}
    \begin{enumerate}
        \item Quelle est l'image de 3 par la fonction $f$ ?
        \item Donner un nombre qui a pour image 5 par la fonction $f$ ?
        \item Donner l'antécédent de 1 par la fonction $f$.
    \end{enumerate}
    \item On considère le programme de calcul suivant :
        \begin{center}
            \begin{mdframed}
            Choisir un nombre\par
            Ajouter 1\par
            Calculer le carré du résultat
            \end{mdframed}
        \end{center}
        \begin{enumerate}
            \item Quel résultat obtient-on en choisissant 1 comme nombre de
            départ ? En choisissant -2 comme nombre de départ ?
            \item On note $x$ le nombre choisi au départ et on appelle $g$ la
            fonction qui a $x$ associe le résultat obtenu avec le programme
            de calcul.\par
            Exprimer $g(x)$ en fonction de $x$.
        \end{enumerate}
        \item La fonction $h$ est définie par $h(x)=2x^2-3$.
        \begin{enumerate}
            \item Quelle est l'image de 3 par la fonction $h$ ?
            \item Quelle est l'image de $-4$ par la fonction $h$ ?
            \item Donner l'antécédent de 5 par la fonction $h$. En existe t-il un autre ?
        \end{enumerate}
\end{enumerate}
\end{document}
